\section{连续嵌入}
\label{sec:sobolev-continuous-embeddings}

一维时，第~\ref{chap:smoothing-weak-derivatives}章已经证明$W^{1,1}$函数有绝对连续代表，所以有界。高维时情况不同：导数可积能把函数本身控制到什么程度，取决于积分指数$p$与维数$d$的关系。本节固定有界Lipschitz区域$\Omega\subset\mathbb R^d$，研究一阶Sobolev范数能控制函数在哪些目标空间中的大小。

若包含关系$W^{1,p}(\Omega)\subset Y$连同估计
\begin{equation}
\lVert f\rVert_Y\leq C\lVert f\rVert_{W^{1,p}}
\label{eq:sobolev-continuous-embedding}
\end{equation}
成立，常数$C$不依赖$f$，就称它为\term{连续嵌入}（continuous embedding）：包含映射是有界线性算子。

本节的两个基本条件各有用途。区域有界保证较小指数的$L^q$范数受较大指数控制：对$q\leq r$，由Hölder不等式，
\[
\lVert f\rVert_q\leq|\Omega|^{1/q-1/r}\lVert f\rVert_r .
\]
区域为Lipschitz，保证第~\ref{sec:sobolev-extension}节的延拓和光滑到边界的逼近可用。

\paragraph{缩放决定指数}
在讨论证明以前，先看目标指数应当是多少。设想全空间上的光滑紧支撑函数满足$\lVert u\rVert_q\leq C\lVert\nabla u\rVert_p$，$C$不依赖$u$。对$\lambda>0$，令$u_\lambda(x)=u(\lambda x)$，换元$y=\lambda x$给出
\[
\lVert u_\lambda\rVert_q=\lambda^{-d/q}\lVert u\rVert_q,\qquad
\lVert\nabla u_\lambda\rVert_p=\lambda^{1-d/p}\lVert\nabla u\rVert_p .
\]

不等式对所有$\lambda$成立，两边$\lambda$的指数必须相同，否则令$\lambda\to0$或$\lambda\to\infty$就会破坏它。于是$-d/q=1-d/p$，即
\begin{equation}
\frac1{p^*}=\frac1p-\frac1d,\qquad p^*=\frac{dp}{d-p}\quad(1\leq p<d).
\label{eq:sobolev-conjugate-exponent}
\end{equation}
称$p^*$为$p$的\term{Sobolev共轭指数}（Sobolev conjugate exponent）。例如$d=3$、$p=2$时$p^*=6$，缩放指数两边都是$\lambda^{-1/2}$。

缩放只说明指数必须是什么，并没有证明不等式成立。它也显示了维数的作用：$d$越大、$p$越小，把函数压缩到小尺度时导数的$L^p$范数增长得越慢，导数就越难察觉尖锐的集中。

\subsection{更高指数的可积性}
\label{sec:sobolev-higher-integrability}

当$p<d$时，我们证明$W^{1,p}$函数属于比$L^p$更好的$L^{p^*}$。

\begin{theorem}[Sobolev嵌入，$p<d$]
\label{thm:sobolev-embedding-subcritical}
设$d\geq2$，$1\leq p<d$，$\Omega\subset\mathbb R^d$为有界Lipschitz区域。存在常数$C$，使所有$f\in W^{1,p}(\Omega)$满足
\begin{equation}
\lVert f\rVert_{L^{p^*}(\Omega)}\leq C\lVert f\rVert_{W^{1,p}(\Omega)} .
\label{eq:sobolev-embedding-subcritical}
\end{equation}
因而对每个$1\leq q\leq p^*$，把左端换成$\lVert f\rVert_q$后估计仍成立，只需把常数乘以$|\Omega|^{1/q-1/p^*}$。
\end{theorem}

证明分四步：先对全空间紧支撑函数证明$p=1$的情形，再借幂函数推到$p>1$，然后用逼近推广到$W^{1,p}(\mathbb R^d)$，最后用延拓回到区域。

\paragraph{第一步：\texorpdfstring{$p=1$}{p=1}的全空间估计}
设$u\in C_c^1(\mathbb R^d)$。沿第$i$个坐标方向，从$-\infty$积分到$x_i$，得到
\[
|u(x)|\leq g_i(\hat x_i):=\int_{\mathbb R}|\partial_iu|\,\mathrm dx_i ,
\]
其中$\hat x_i$表示去掉第$i$个坐标后的变量，右端不依赖$x_i$。这里每个坐标方向都给了一个上界。

先看$d=2$。此时$|u(x)|^2\leq g_1(x_2)g_2(x_1)$，右端是两个单变量函数之积，对$x_1,x_2$积分后分离：
\[
\int_{\mathbb R^2}|u|^2\,\mathrm dx\leq\Bigl(\int g_1\Bigr)\Bigl(\int g_2\Bigr)=\lVert\partial_1u\rVert_1\lVert\partial_2u\rVert_1 .
\]
于是$\lVert u\rVert_2\leq\lVert\nabla u\rVert_1$，这正是$d=2$、$p=1$、$p^*=2$的情形。

一般维数需要一个把$d$个“各少一个变量”的函数之积分离开的不等式。

\begin{lemma}[分离变量的乘积估计]
\label{lem:sobolev-loomis-whitney}
设$d\geq2$，$h_1,\ldots,h_d$是$\mathbb R^{d-1}$上的非负可测函数，$h_i$的自变量是$\hat x_i$。则
\begin{equation}
\int_{\mathbb R^d}\prod_{i=1}^dh_i(\hat x_i)\,\mathrm dx\leq\prod_{i=1}^d\lVert h_i\rVert_{L^{d-1}(\mathbb R^{d-1})} .
\label{eq:sobolev-loomis-whitney}
\end{equation}
\end{lemma}

\begin{proof}
对维数归纳。$d=2$时，积分直接分离：
\[
\int_{\mathbb R^2}h_1(x_2)\,h_2(x_1)\,\mathrm dx=\lVert h_1\rVert_1\lVert h_2\rVert_1 .
\]

设维数$d\geq2$时成立，考虑$\mathbb R^{d+1}$上的$h_1,\ldots,h_{d+1}$。记最后一个变量为$s$，其余变量为$x'\in\mathbb R^d$，则$h_{d+1}=h_{d+1}(x')$不依赖$s$。固定$s$，对$x'$积分，用指数$d$和$d/(d-1)$的Hölder不等式，得
\[
\int\prod_{i=1}^{d+1}h_i\,\mathrm dx'
\leq\lVert h_{d+1}\rVert_{L^d}\Bigl(\int\prod_{i=1}^dh_i^{d/(d-1)}\,\mathrm dx'\Bigr)^{(d-1)/d}.
\]

对括号内的积分用归纳假设，函数换成$h_i^{d/(d-1)}$。记$A_i(s)=\int h_i(\cdot,s)^d$，积分对$h_i$的其余$d-1$个变量进行。归纳假设给出
\[
\Bigl(\int\prod_{i=1}^dh_i^{d/(d-1)}\,\mathrm dx'\Bigr)^{(d-1)/d}\leq\prod_{i=1}^dA_i(s)^{1/d}.
\]

最后对$s$积分，用$d$个指数都等于$d$的推广Hölder不等式：
\[
\int\prod_{i=1}^dA_i(s)^{1/d}\,\mathrm ds\leq\prod_{i=1}^d\Bigl(\int A_i(s)\,\mathrm ds\Bigr)^{1/d}=\prod_{i=1}^d\lVert h_i\rVert_{L^d(\mathbb R^d)} .
\]
三式合起来，就是维数$d+1$时的结论。
\end{proof}

回到估计。把$d$个坐标方向的界相乘，再开$d-1$次方，得到
\[
|u(x)|^{d/(d-1)}\leq\prod_{i=1}^dg_i(\hat x_i)^{1/(d-1)} .
\]
对$h_i=g_i^{1/(d-1)}$用引理~\ref{lem:sobolev-loomis-whitney}。按$g_i$的定义，
\[
\lVert h_i\rVert_{d-1}^{d-1}=\int g_i\,\mathrm d\hat x_i=\lVert\partial_iu\rVert_1,
\]
所以
\[
\int_{\mathbb R^d}|u|^{d/(d-1)}\,\mathrm dx\leq\prod_{i=1}^d\lVert\partial_iu\rVert_1^{1/(d-1)} .
\]
两边开$(d-1)/d$次方，得到
\begin{equation}
\lVert u\rVert_{d/(d-1)}\leq\prod_{i=1}^d\lVert\partial_iu\rVert_1^{1/d}\leq\lVert\nabla u\rVert_1 .
\label{eq:sobolev-gagliardo-nirenberg}
\end{equation}

\paragraph{第二步：\texorpdfstring{$1<p<d$}{1<p<d}}
对$u\in C_c^1(\mathbb R^d)$和$\gamma>1$，$|u|^\gamma$仍是紧支撑的连续可微函数，梯度满足
\[
\bigl|\nabla|u|^\gamma\bigr|\leq\gamma|u|^{\gamma-1}|\nabla u| .
\]
把式~\eqref{eq:sobolev-gagliardo-nirenberg}用于$|u|^\gamma$，再对右端用Hölder不等式，
\[
\lVert u\rVert_{\gamma d/(d-1)}^\gamma
\leq\gamma\int|u|^{\gamma-1}|\nabla u|\,\mathrm dx
\leq\gamma\lVert u\rVert_{(\gamma-1)p'}^{\gamma-1}\lVert\nabla u\rVert_p .
\]

选择$\gamma=p(d-1)/(d-p)$，让两侧出现同一个范数。逐项核对：
\[
\frac{\gamma d}{d-1}=\frac{dp}{d-p}=p^*,\qquad
(\gamma-1)p'=\frac{d(p-1)}{d-p}\cdot\frac p{p-1}=p^*.
\]
$u$紧支撑，$\lVert u\rVert_{p^*}$有限，若不为零就可以约去$\gamma-1$次方，得到
\begin{equation}
\lVert u\rVert_{p^*}\leq\gamma\lVert\nabla u\rVert_p,\qquad\gamma=\frac{p(d-1)}{d-p}.
\label{eq:sobolev-gns-p}
\end{equation}
$p=1$时$\gamma=1$，与第一步一致。$d=3$、$p=2$时$\gamma=4$，即$\lVert u\rVert_6\leq4\lVert\nabla u\rVert_2$。

\paragraph{第三步：推广到全空间Sobolev函数}
由第~\ref{sec:sobolev-compact-closure}小节，$C_c^\infty(\mathbb R^d)$在$W^{1,p}(\mathbb R^d)$中稠密。对$u\in W^{1,p}(\mathbb R^d)$取$u_n\in C_c^\infty(\mathbb R^d)$在$W^{1,p}$中趋于$u$，再取子列几乎处处收敛。由Fatou引理与式~\eqref{eq:sobolev-gns-p}，
\[
\lVert u\rVert_{p^*}\leq\liminf_n\lVert u_n\rVert_{p^*}\leq\gamma\lim_n\lVert\nabla u_n\rVert_p=\gamma\lVert\nabla u\rVert_p .
\]

\paragraph{第四步：回到区域}
对$f\in W^{1,p}(\Omega)$，用定理~\ref{thm:sobolev-extension}的延拓$Ef$。比较$\ell^1$与$\ell^p$范数，把梯度的范数换成Sobolev范数，会带来因子$d^{1-1/p}$，于是
\[
\lVert f\rVert_{L^{p^*}(\Omega)}\leq\lVert Ef\rVert_{L^{p^*}(\mathbb R^d)}
\leq\gamma\,d^{1-1/p}\,\lVert E\rVert\,\lVert f\rVert_{W^{1,p}(\Omega)},
\]
其中$\lVert E\rVert$是延拓算子的范数。较小指数由区域有限测度得到，这就证明了定理~\ref{thm:sobolev-embedding-subcritical}。

常数的来源由此清楚：$\gamma$只依赖$d,p$，$\lVert E\rVert$依赖区域与$p$，较小指数还多出$|\Omega|$的幂。它们都不依赖被估计的函数，这正是连续嵌入所要求的统一性。

\begin{example}[幂函数与指数的精确性]
\label{ex:sobolev-power-function}
在单位球$B\subset\mathbb R^d$中取$u(x)=|x|^{-a}$，$a>0$。它在原点之外光滑，$|\nabla u|=a|x|^{-a-1}$。若$a+1<d$，挖去半径为$r$的小球后分部积分，小球面上的边界项不超过
\[
\sup|\varphi|\cdot r^{-a}\cdot|S^{d-1}|\,r^{d-1},
\]
随$r\to0$趋于零，所以经典梯度就是弱梯度。

按球坐标，$|x|^{-b}$在$B$上可积当且仅当$b<d$，因此
\[
u\in L^q(B)\iff aq<d,\qquad u\in W^{1,p}(B)\iff(a+1)p<d .
\]
后一条件即$a<d/p-1=d/p^*$，它恰好是$u\in L^{p^*}(B)$的条件。

对任何$q>p^*$，取$a$满足$d/q\leq a<d/p-1$，便得到$W^{1,p}(B)$中不属于$L^q(B)$的函数，所以$p^*$不能再提高。取$d=3$、$p=2$、$p^*=6$，再取$q=8$与$a=0.4$，逐项核对：
\[
(a+1)p=2.8<3,\qquad 6a=2.4<3,\qquad 8a=3.2\geq3 .
\]
所以这个$u$属于$W^{1,2}(B)$和$L^6(B)$，不属于$L^8(B)$，并且在原点附近无界。可见$p<d$时得到的是更高的可积性，而不是有界性或连续性。
\end{example}

\subsection{临界指数下的可积性}
\label{sec:sobolev-critical-integrability}

当$p\to d$时，$p^*=dp/(d-p)\to\infty$。这提示我们考察$p=d\geq2$时能得到什么。

\begin{theorem}[临界指数的嵌入]
\label{thm:sobolev-embedding-critical}
设$d\geq2$，$\Omega\subset\mathbb R^d$为有界Lipschitz区域。对每个有限的$q\geq1$，存在常数$C_q$，使所有$f\in W^{1,d}(\Omega)$满足
\begin{equation}
\lVert f\rVert_q\leq C_q\lVert f\rVert_{W^{1,d}} .
\label{eq:sobolev-embedding-critical}
\end{equation}
\end{theorem}

\begin{proof}
给定$q$，取$1\leq s<d$足够接近$d$，使$s^*=ds/(d-s)\geq q$。区域有界，由Hölder不等式，
\[
\lVert g\rVert_s\leq|\Omega|^{1/s-1/d}\lVert g\rVert_d
\]
对$f$及其各一阶弱导数成立，所以$W^{1,s}$范数受$W^{1,d}$范数控制。再用定理~\ref{thm:sobolev-embedding-subcritical}于指数$s$即得。
\end{proof}

事实上，$W^{1,d}$函数一般无界。在单位球$B\subset\mathbb R^d$中取$u(x)=\ln\ln(e/|x|)$。记$r=|x|$，$t=\ln(e/r)\geq1$，则$|\nabla u|=1/(rt)$。按球坐标，
\[
\begin{aligned}
\int_B|\nabla u|^d\,\mathrm dx
&=|S^{d-1}|\int_0^1\frac{r^{d-1}}{r^d(\ln(e/r))^d}\,\mathrm dr\\
&=|S^{d-1}|\int_1^\infty t^{-d}\,\mathrm dt=\frac{|S^{d-1}|}{d-1},
\end{aligned}
\]
$d\geq2$时有限。这里用了换元$t=\ln(e/r)$，$\mathrm dt=-\mathrm dr/r$。

$u$本身只以$\ln\ln$的速度增长，属于每个$L^q(B)$。小球面上的边界项不超过
\[
\sup|\varphi|\cdot u(r)\,|S^{d-1}|\,r^{d-1}\to0,
\]
所以经典梯度是弱梯度。于是$u\in W^{1,d}(B)$，却在原点趋于无穷。

这说明不能把式~\eqref{eq:sobolev-conjugate-exponent}在$p=d$处“代入”成$p^*=\infty$。常数$C_q$也不可能关于$q$一致：若$C_q\leq C$对所有$q$成立，由于在有限测度区域上$q\to\infty$时$\lVert f\rVert_q$趋于$\lVert f\rVert_\infty$，就得到
\[
\lVert f\rVert_\infty\leq C\lVert f\rVert_{W^{1,d}},
\]
与这个反例矛盾。临界情形处在两种结论之间：可积性对每个有限指数都成立，有界性却失去了。

\subsection{连续代表}
\label{sec:sobolev-continuous-representatives}

当$p>d$时，导数的控制足以给出逐点的振荡估计。机制是比较一点的值与它附近的平均：引理~\ref{lem:sobolev-convex-average}把这个差写成梯度乘以核$|x-z|^{1-d}$的积分，而这个核的$p'$次幂恰好在$p>d$时局部可积。

\begin{theorem}[Hölder连续代表]
\label{thm:sobolev-morrey}
设$d<p<\infty$，$\Omega\subset\mathbb R^d$为有界Lipschitz区域，$\beta=1-d/p$。每个$f\in W^{1,p}(\Omega)$有一个在$\overline\Omega$上连续的代表$\tilde f$，并且存在不依赖$f$的常数$C$，使
\begin{equation}
\sup_{x\in\overline\Omega}|\tilde f(x)|
+\sup_{x\neq y}\frac{|\tilde f(x)-\tilde f(y)|}{|x-y|^\beta}
\leq C\lVert f\rVert_{W^{1,p}(\Omega)} .
\label{eq:sobolev-morrey}
\end{equation}
\end{theorem}

\begin{proof}
先设$u\in C_c^1(\mathbb R^d)$。取两点$x\neq y$，$r=|x-y|$，令$W=B(x,r)\cap B(y,r)$。$W$是凸开集，直径不超过$2r$，并且包含以中点为心、半径$r/2$的球，所以$|W|\geq2^{-d}|B(0,1)|\,r^d$。

$x,y\in\overline W$，由引理~\ref{lem:sobolev-convex-average}与$W\subset B(x,r)$，存在只依赖$d$的$C_d$，使
\[
\begin{aligned}
|u(x)-u_W|&\leq C_d\int_{B(x,r)}\frac{|\nabla u(z)|}{|x-z|^{d-1}}\,\mathrm dz\\
&\leq C_d\lVert\nabla u\rVert_p\Bigl(\int_{B(x,r)}|x-z|^{(1-d)p'}\,\mathrm dz\Bigr)^{1/p'} .
\end{aligned}
\]
按球坐标，最后的积分等于
\[
\frac{|S^{d-1}|\,r^{d-(d-1)p'}}{d-(d-1)p'},
\]
它有限当且仅当$(d-1)p'<d$，即$p>d$。其$1/p'$次方是常数乘$r^{1-d/p}$。

对$y$作同样估计，两式相加，得到
\[
|u(x)-u(y)|\leq C|x-y|^\beta\lVert\nabla u\rVert_p .
\]
取$W=B(x,1)$，又有
\[
|u(x)|\leq|u_W|+C\lVert\nabla u\rVert_p\leq|B(0,1)|^{-1/p}\lVert u\rVert_p+C\lVert\nabla u\rVert_p .
\]

对$f\in W^{1,p}(\Omega)$，令$u_n=(Ef)*\rho_{1/n}\in C_c^\infty(\mathbb R^d)$，它们在$W^{1,p}(\mathbb R^d)$中趋于$Ef$。把上面的上确界估计用于$u_n-u_m$，可知$(u_n)$一致收敛到一个连续函数$\tilde u$。$u_n$也在$L^p$中趋于$Ef$，所以$\tilde u=Ef$几乎处处，$\tilde f=\tilde u|_{\overline\Omega}$是$f$的代表。Hölder估计和上确界估计在一致极限下保持，再由延拓算子有界即得结论。
\end{proof}

满足上述估计的函数称为$\beta$阶\term{Hölder连续}（Hölder continuous）。连续代表在等价类中是唯一的：两个连续函数若几乎处处相等，就处处相等。所以$p>d$时，点值$f\mapsto\tilde f(x)$是$W^{1,p}(\Omega)$上的有界线性泛函，这与$p\leq d$时一般无法读取点值形成对照。

指数$\beta$不能整体提高。在单位球中取$u(x)=|x|^b$，$0<b<1$，它的梯度长度是$b|x|^{b-1}$，属于$L^p$当且仅当$(1-b)p<d$，即$b>1-d/p$。这样的$b$可以任意接近$1-d/p$，而$u$在原点处恰好只有$b$阶Hölder连续。例如$d=2$、$p=4$时$\beta=1/2$，取$b=0.6$，有$(1-b)p=1.6<2$。

$p=\infty$时，指数$\beta$趋于$1$，相应的结论是Lipschitz代表。先在凸区域上说明。设$\Omega$凸，$f\in W^{1,\infty}(\Omega)$。由式~\eqref{eq:sobolev-mollifier-commute}，在$\Omega_\varepsilon$上$f_\varepsilon$的梯度是$\nabla f$的磨光，其长度不超过$\lVert\nabla f\rVert_\infty$。对$x,y\in\Omega$，线段$[x,y]$是$\Omega$中的紧集，$\varepsilon$足够小时落在$\Omega_\varepsilon$中，所以
\[
|f_\varepsilon(x)-f_\varepsilon(y)|\leq\lVert\nabla f\rVert_\infty|x-y| .
\]
取$\varepsilon_n\to0$使$f_{\varepsilon_n}\to f$几乎处处，极限在一个满测度集合上满足同一Lipschitz估计，再唯一地连续延伸到整个$\Omega$，得到$\tilde f$。

对一般有界Lipschitz区域，同样的论证在每个含于区域的球上成立。区域只有有限个连通分支，不同分支之间距离为正，用$\lVert f\rVert_\infty$即可控制；同一分支中任意两点可以用长度不超过$C|x-y|$的折线相连（这一几何事实这里不证明），于是代表在整个区域上Lipschitz，常数为$C\lVert f\rVert_{W^{1,\infty}}$。

\paragraph{一维的情形}
$d=1$时，有界区间$(a,b)$上的$W^{1,1}$函数已经有绝对连续代表，与端点估计同样的平均论证给出
\[
|\tilde f(y)|\leq\frac1{b-a}\int_a^b|f|\,\mathrm dx+\int_a^b|f'|\,\mathrm dx .
\]
所以一维时$p=d=1$不仅有每个有限指数的可积性，还有有界的连续代表。$p>1$时，由
\[
|\tilde f(x)-\tilde f(y)|\leq\int_x^y|f'|\,\mathrm dt\leq|x-y|^{1-1/p}\lVert f'\rVert_p,
\]
代表是$1-1/p$阶Hölder连续的，与$\beta=1-d/p$一致。

$d\geq2$的临界反例$\ln\ln(e/|x|)$依赖函数能够绕开一点集中增长，而一维时从一点到另一点只有一条路径，变化必须全部经过它。因此不能把$d\geq2$的临界结论照搬到$d=1$。

\paragraph{高阶空间}
同样的指数比较可以推广到$W^{k,p}$：每多一阶导数，缩放中多出一个因子$\lambda$。若$kp<d$，则$W^{k,p}(\Omega)\subset L^q(\Omega)$，$1/q=1/p-k/d$；若$kp>d$，则函数有连续代表。特别地，$k>d/2$时$H^k(\Omega)$的元素有连续代表，且$\sup|\tilde f|\leq C\lVert f\rVert_{H^k}$。一般情形的证明要把前面的一阶论证反复叠加，篇幅较长，读者可查阅相关专著。

坐标的图景可以解释条件$k>d/2$。在周期函数上，$H^k$的权约为$(1+|\xi|^2)^k$，$\xi\in\mathbb Z^d$是频率向量。和第~\ref{sec:sobolev-endpoint-values}小节读取点值的计算一样，点值有界要求
\[
\sum_{\xi\in\mathbb Z^d}\frac1{(1+|\xi|^2)^k}<\infty .
\]
半径约为$R$的球壳中大约有$R^{d-1}$个频率点，每个贡献$R^{-2k}$，这个和收敛当且仅当$2k>d$。维数升高时频率点变多，权必须增长得更快才能把点值压成有限，这就是一维$H^1$能读点值而二维$H^1$不能的原因。它使每个点值读取都是$H^k$上的有界线性泛函，第~\ref{chap:rkhs}章将以这类空间为例研究连续点值。

\paragraph{四种情形的对照}
表~\ref{tab:sobolev-embedding-cases}汇总了一阶空间在维数$d\geq2$时的结论，最后一列是说明结论不能再改进的例子。

\begin{BookInlineTable}
  \centering
  \caption{$W^{1,p}$在有界Lipschitz区域上的连续嵌入（$d\geq2$）。}
  \label{tab:sobolev-embedding-cases}
  \begin{tabular}{lll}
    \toprule
    指数 & 能控制到 & 例子\\
    \midrule
    $1\leq p<d$ & $L^{p^*}$，$p^*=dp/(d-p)$ & $|x|^{-a}$\\
    $p=d$ & 每个有限的$L^q$ & $\ln\ln(e/|x|)$\\
    $d<p<\infty$ & $1-d/p$阶Hölder代表 & $|x|^b$，$b$略大于$1-d/p$\\
    $p=\infty$ & Lipschitz代表 & $|x-x_0|$\\
    \bottomrule
  \end{tabular}
\end{BookInlineTable}
